\section{Introduction}
\label{intro:section}

Zeros in an analytic family can be counted in several ways. The argument
principle counts complex zeros with multiplicity. A real heat evolution
can turn a conjugate pair into two real zeros, and a real spacetime index
records that event. A complex-time discriminant counts collisions with
yet another multiplicity. On a branched cover, the same zeros determine
the topology, while a chosen differential supplies an analytic energy.
The usefulness of a boundary balance depends on distinguishing these
quantities and retaining the complete boundary terms.

This paper develops a self-contained finite-domain framework linking
these descriptions. Its main analytic identity is particularly simple.
If $f$ is holomorphic near a bounded domain $\Omega$, has no boundary
zero, and has $N$ interior zeros with multiplicity, then
\begin{equation}\label{intro:main}
 2\int_{\Omega\cap\{|f|<\varepsilon\}}
       \frac{|f'|^2}{|f|}\,dA=4\pi\varepsilon N,
 \qquad0<\varepsilon<\min_{\partial\Omega}|f|.
\end{equation}
There is no simplicity or separation hypothesis. The integral remains
locally finite through every finite multiplicity. On the normalized
cover $w^2=f(z)$ it is the energy of the exact differential $2\,dw$
over the low-value region. More generally, the energy of $q\,dw$ on
$w^q=f(z)$ is exactly $\pi q^2\varepsilon^{2/q}N$ on the corresponding
region. These are applications of the holomorphic degree/area formula;
the point is their quantitative formulation and their interaction with
transport, rather than a new version of the argument principle.

The localized energy has a compact target-space potential. This gives a
regularized divisor measure of exact total mass $N$, an approximation
allowance proportional to the field error divided by the value cutoff,
and a quantitative weak-limit bias near roots of arbitrary multiplicity.
A family of power potentials connects the double-cover case to higher
root covers. We also compare the exact-mass regularization with a smooth
logarithmic kernel whose smaller Lipschitz constant comes with an
explicit boundary bias. These statements specify the information needed
to use an energy integral as a count; they do not infer a count from an
uncontrolled normalization. Harmonic reproduction removes the linear
localization error for $C^2$ tests. With complete disk, boundary and
analytic-unit data, a direct contour estimate gives weighted-divisor
stability of order $\eta^{\min(1,2/m)}$ at an $m$-fold zero;
monomial and translation examples attain the exponent.

For the backward heat equation
\begin{equation}\label{intro:heat}
 H_t=-H_{zz},
\end{equation}
the complex multiplicity count is conserved when roots do not cross the
spatial boundary. Real roots obey a different complete balance: an
interior spatial multiplicity $m$ contributes $2\lfloor m/2\rfloor$
new real roots, with signed wall crossings added to the cap difference.
The complex discriminant instead has local vanishing order
$\binom m2$. We prove these weights from the local Hermite expansion and
give the complete cap conventions. The normalized double cover has a
singular-fiber Euler-characteristic excess
$2\sum\lfloor m/2\rfloor$ over a regular fiber with the same total
complex count. These formulas show precisely when transported simple
zeros cancel and what remains after that cancellation.

Root geometry provides another analytic interface. Inverse-square pair
distances form a graph Laplacian, the Hessian of logarithmic Vandermonde
energy. We derive its complete heat evolution, retaining the holomorphic
factor containing the omitted roots. A full polynomial heat flow has a
strongly convex rescaled Hermite energy. For a selected cluster of an
entire function, the exterior factor acts as a force. We give an explicit
matrix comparison and a self-contained large-time application to the
complete theta source: blocks ending by rank $K=o(L\sqrt\kappa)$
converge after scaling to their clock matrices with a quantitative error
uniform in the block size. A fixed-population cumulative cancellation
formula explains why the exterior force cannot be discarded.

Pair correlation is a possible additional input, but it uses a different
limiting topology. We compute the clock and sine-kernel inverse-square
moments, prove a weighted transfer theorem with its small-gap and tail
hypotheses, and give an explicit finite-population error threshold that
would exclude repeated points. Finally, the physical projective cap
formula shows that the root-derivative phase is additional information:
regular simple roots alone do not make the cap integral vanish. Its two
imaginary-time faces are evaluated with their orientations.

\paragraph{Analytic scope and relation to earlier work.}
The basic ingredients are classical: the argument principle, the
holomorphic change-of-variables formula, Green's identity,
Riemann--Hurwitz, Hermite polynomials, and logarithmic root interactions.
The real and complex transport identities consolidate earlier work in
the accompanying research repository. We provide their proofs here so
that the present paper does not require those internal notes. The
regularization family, explicit stability and localization estimates,
matrix bounds, complete-source concentration estimate and boundary
phase formulas are the quantitative statements to be assessed on their
own hypotheses. No claim of mathematical priority is made.

The theta field uses the standard normalization in
Rodgers--Tao~\cite{RT}. Our application concerns large positive heat
time and the specified growing spatial windows, not a conjectured
high-height pair law. Cumulative dynamical identities retain fixed populations.
None of the main analytic results assumes RH or proves RH or endpoint
simplicity. The only numerical input appears in
Appendix~\ref{app:numerics}, as an illustration using a previously
certified finite disk. Removing that appendix changes none of the
proofs or analytic conclusions.

\paragraph{Conventions.}
Spatial area is $dA=dx\,dy$, and the energy of a differential $\eta$ on
a Riemann surface is $(i/2)\int\eta\wedge\bar\eta$. Complex time is
$\tau=t+is$. Ordered pair measures omit only the label diagonal;
coincident locations with different labels remain counted. Every
boundary orientation is stated when used. The notation $O_K$ allows
constants depending on a fixed finite rank bound $K$; it does not assert
uniformity as the population grows.
